\chapter{糖：燃料、材料和身份标签}

\begin{kaichang}
请你去凑四样东西：一块方糖、一勺白米饭、一张纸巾（或者一小团棉花），还有体检报告上印着血型的那一栏——比如“A 型”“O 型”。

现在猜一猜：这四样东西之间，有什么共同的“原料”？方糖和米饭都能吃，纸巾可不能吃；血型更是看不见摸不着，得抽血才查得出来。它们凭什么算一家人？

把你的猜测写在纸上。这一章结束时你会发现：它们的主要成分竟然是\textbf{同一种小分子}——葡萄糖。区别只在于搭法：接成什么链、用什么角度接头、挂在什么东西表面。同一种积木，搭出了燃料、建筑材料和身份标签三种命运。
\end{kaichang}

\section{先只信眼睛和舌头}

现象层，先不急着讲道理，只看能看、能尝、能测的。

方糖：白色晶体，丢进水里搅一搅就消失，尝起来甜。放进烤箱加热到一百六七十摄氏度，它会融化、变黄、变褐，飘出焦糖香——颜色和气味都变了，说明有化学反应发生。

白米饭：刚入口几乎不甜。但如果你忍住不咽，慢慢嚼上一分钟，会尝出淡淡的甜味。饭还是那口饭，甜味是从哪里来的？

纸巾和棉花：不甜，泡在水里一整天也不溶解，用牙咬也咬不烂。可它们点得着，烧完的灰烬和糖烧焦后的残渣一样黑。

血糖仪：扎一滴血，读出一个数字。空腹时健康人的血里大约每升含 1 克葡萄糖。注意，仪器报的是“葡萄糖”这一种分子，不是泛指一切糖。

血型：试纸上的抗体遇到你的血，凝集或不凝集。这一“认”一“不认”，认的到底是什么？

五种现象看似风马牛不相及。把它们缝到一起的线头，在分子层面。

\section{六碳小环：葡萄糖的真面目}

把镜头推进到分子。方糖、米饭、纸巾里含量最高的“砖块”是同一个分子：\textbf{葡萄糖}——6 个碳、12 个氢、6 个氧。碳骨架排成一条六个碳的短链，链上挂满羟基（\chem{-OH}，第 40 章讲过，它亲水），链头还有一个醛基。浑身羟基，这就是糖爱溶于水的原因：水分子爱跟羟基拉氢键（第 19 章），一拥而上把糖分子“抬”进水里。

但在水溶液里，这条链很少以伸直的样子存在。链头的醛基是个活泼的接口，而第 5 个碳上的羟基正好够得着它——于是分子咬住自己的尾巴：5 号碳的羟基与 1 号碳的醛基反应，首尾相连，蜷成一个六元环，环上 5 个碳加 1 个氧。在水里，$99\%$ 以上的葡萄糖都是环，伸直的链不到千分之一。

\begin{center}
\begin{tikzpicture}[scale=1.0]
\coordinate (O) at (0.9,0.9);
\coordinate (C1) at (1.6,0);
\coordinate (C2) at (0.9,-0.9);
\coordinate (C3) at (-0.9,-0.9);
\coordinate (C4) at (-1.6,0);
\coordinate (C5) at (-0.9,0.9);
\draw[thick] (O)--(C1)--(C2)--(C3)--(C4)--(C5)--cycle;
\node at (0.95,0.95) {\chem{O}};
\node[below] at (1.75,-0.15) {\small 1};
\node[below] at (0.9,-1.05) {\small 2};
\node[below] at (-0.9,-1.05) {\small 3};
\node[below] at (-1.85,-0.15) {\small 4};
\node[above] at (-0.9,1.05) {\small 5};
\draw[thick] (C1) -- (1.6,-0.7) node[below] {\small \chem{OH}};
\draw[thick] (C5) -- (-0.9,1.7) node[above] {\small \chem{CH_2OH}};
\node at (0,-1.9) {\small 1 号碳：接头碳，羟基朝下是 $\alpha$、朝上是 $\beta$};
\end{tikzpicture}
\end{center}

这是人为的表示法（化学家叫它哈沃斯投影）：环其实不平，羟基的位置也只是相对关系的记号，颜色和大小都不代表真实外观。

关环的时候出了一个选择：1 号碳上新长出来的那个羟基，可以朝环平面的下方，也可以朝上方。朝下叫 $\alpha$（阿尔法），朝上叫 $\beta$（贝塔）——纯粹是人为约定的记号。就这一个羟基的朝向，让葡萄糖分成了两种版本。

再往深看一层：第 42 章讲过手性，一个碳连着四个不同的基团时，空间排布不同就是不同的分子。葡萄糖身上有四个这样的手性中心，理论上有 16 种立体异构体——葡萄糖只是其中之一，半乳糖、甘露糖都是它的“异构兄弟”，成分一模一样，差别只在某个碳上羟基的朝向。而地球上的生命几乎只用其中一种版本。为什么生命如此挑剔？证据故事一节再讲，而且答案里有一块至今空着。

“羟基朝上还是朝下”这种看似吹毛求疵的细节，请先记住。后面你会看到，纸张能不能吃、牛奶喝了会不会拉肚子，全由这种朝向决定。

\section{一节一节接起来：双糖和多糖}

单个的糖叫\textbf{单糖}：葡萄糖、果糖、半乳糖都是。两个单糖可以“牵手”：一边出羟基，另一边出羟基上的氢，两处羟基之间脱掉一分子水，连上一座氧桥——这座 \chem{C-O-C} 的桥叫\textbf{糖苷键}。这正是第 43 章缩合聚合的路数：每接一节，付一笔“水费”。

两个单糖接起来叫\textbf{双糖}。厨房里最常见的三种甜味都来自双糖：

\begin{itemize}[itemsep=2pt]
\item 蔗糖（方糖、白糖、红糖的主要成分）：葡萄糖 $+$ 果糖；
\item 乳糖（牛奶里的糖）：葡萄糖 $+$ 半乳糖；
\item 麦芽糖（糖稀、麦芽）：葡萄糖 $+$ 葡萄糖。
\end{itemize}

成千上万个单糖接起来，就是\textbf{多糖}。这里发生了一件值得停下来盯着看的事：淀粉、糖原、纤维素，三种多糖的单体\textbf{一模一样，全是葡萄糖}。可淀粉是米饭面粉的主料，糖原是藏在肝脏和肌肉里的储备粮，纤维素是棉花和纸——三种东西的硬度、溶解性、能不能吃，天差地别。同样的砖，凭什么盖出完全不同的房子？

差别全在接头的朝向和位置上。

淀粉用 $\alpha$ 型接头：一个葡萄糖 1 号碳的 $\alpha$ 羟基，与下一个葡萄糖的 4 号碳相连（记作 $\alpha$-1,4）。这个接头带一个“拐”，链天然地弯成螺旋，像弹簧一样盘起来，松松软软，水分子钻得进去。支链淀粉和糖原还每隔一二十个葡萄糖，就在 6 号碳上岔出一根侧枝（$\alpha$-1,6 接头），整个分子像一棵灌木。糖原分支最密——这个细节章末有用。

纤维素用 $\beta$ 型接头（$\beta$-1,4）。别小看这一个朝向的差别：$\beta$ 接头不带拐，链是笔直的，而且每个葡萄糖相对前一个要翻个身。笔直的链并排躺在一起，相邻链上的羟基正好面对面，拉起密密麻麻的氢键，把几十条链捆成一束结实的纤维——就像一把散筷子捆成了筷子束。水分子挤不进这张氢键网，所以纸泡水不烂；纤维的强度高到可以做衣服的布料。

\begin{center}
\begin{tikzpicture}[scale=0.9]
\foreach \x/\y in {0/0, 1.4/-0.4, 2.8/0, 4.2/-0.4} \draw (\x,\y) circle (0.3);
\draw (0.3,-0.1) -- (1.1,-0.3); \draw (1.7,-0.3) -- (2.5,-0.1);
\draw (3.1,-0.1) -- (3.9,-0.3);
\node at (5.6,-0.2) {\small $\alpha$-1,4（淀粉）：链打弯};
\foreach \x in {0, 1.4, 2.8, 4.2} \draw (\x,-1.8) circle (0.3);
\draw (0.3,-1.8) -- (1.1,-1.8); \draw (1.7,-1.8) -- (2.5,-1.8); \draw (3.1,-1.8) -- (3.9,-1.8);
\draw[dashed] (0,-2.15) -- (0,-2.6); \draw[dashed] (1.4,-2.15) -- (1.4,-2.6);
\draw[dashed] (2.8,-2.15) -- (2.8,-2.6); \draw[dashed] (4.2,-2.15) -- (4.2,-2.6);
\node at (5.6,-1.8) {\small $\beta$-1,4（纤维素）：链绷直};
\node at (2.1,-2.9) {\small 虚线：相邻链之间的氢键};
\end{tikzpicture}
\end{center}

圆圈代表一个个葡萄糖环，连线代表糖苷键——同样是人为的表示法。你看，积木完全相同，只是接头的立体角度不同，一种成了入口的粮食，一种成了挡风的材料。\textbf{结构决定性质}——这句从第 3 章一直说到现在的话，在糖身上演示得再清楚不过。

\section{糖苷键的两笔账：能量和速率}

糖苷键怎么拆？逆着来：把水加回去，糖苷键断开，两端的羟基物归原主。这叫\textbf{水解}。你从米饭里获得能量，第一步就是把淀粉水解成麦芽糖、再水解成葡萄糖。

水解在能量上是“顺坡”的：把多糖拆小，体系的总能量降低（第 27 章）。可是——一碗淀粉汤放在桌上，放一年也不会自己变甜；一张纸泡在水里，十年也化不成糖水。能量上允许，现实中不动。为什么？第 28 章的速率规则：拆糖苷键要先翻过一座活化能的小山，常温下能翻过去的分子寥寥无几。

要拆得快，得请催化剂把山削低。你唾液里的淀粉酶就是这种催化剂——这就是嚼米饭变甜的秘密：酶在几分钟内把淀粉水解出麦芽糖，甜味冒了出来。酶的底细第 49 章细讲，这里只记住两条：酶只改变反应路径、降低活化能，它不提供能量，也不改变反应最终能走多远；而且酶挑接头——淀粉酶只认 $\alpha$ 接头，面对纤维素的 $\beta$ 接头，它无从下手。

这就是“人不能消化纤维素”的真正原因：不是纤维素“有毒”或者“太硬”，而是我们的身体没有配那把开 $\beta$ 接头的钥匙。牛和羊自己也没有——它们靠的是瘤胃里养的微生物，微生物产的酶会拆 $\beta$ 接头。所以牛吃草长肉，本质上是微生物先吃草，牛再分享微生物的劳动成果。白蚁吃木头，用的是同一套路。

\begin{qiaoliang}
数一数你全身有多少葡萄糖。正常空腹血糖浓度大约是每升血 \ud{1}{g}（医学上报的是毫摩尔每升，空腹约 $3.9$～$6.1\ \mathrm{mmol/L}$，换算过来就在每升 1 克上下）。成年人全身血液约 5 升，乘一下：

\[1\ \mathrm{g/L}\times 5\ \mathrm{L}=5\ \mathrm{g}\]

你全身血液里的葡萄糖，总共只有约 5 克——一茶匙那么多。这 5 克全部烧掉能放多少能量？葡萄糖完全氧化每摩尔（180 克）放出约 \ud{2870}{kJ}，所以：

\[\frac{5}{180}\times 2870\approx 80\ \mathrm{kJ}\]

而一个静坐着的人，一天的基础代谢就要约 \ud{7000}{kJ}——血糖里这 80 kJ，只够烧十几分钟。结论很清楚：血糖这点存货必须随用随补（靠肝糖原放粮），却又不能存太多（血糖过高会惹出后面的麻烦）。是谁在精确维持这个“一茶匙”的平衡？胰岛素等激素的调节故事，后面的章节会讲。
\end{qiaoliang}

\section{轮到符号登场}

忍着看到这里，符号可以出来了。葡萄糖的分子式：

\[\chem{C_6H_{12}O_6}\]

注意它只报“成分”，不报结构——果糖、半乳糖的成分一模一样，也是 \chem{C_6H_{12}O_6}。分子式看不出羟基的朝向，这就是为什么本章花了这么多篇幅讲环和接头：化学式是压缩包，结构才是正文。

蔗糖的水解：

\[\chem{C_{12}H_{22}O_{11}} + \chem{H_2O} \rightarrow \chem{C_6H_{12}O_6} + \chem{C_6H_{12}O_6}\]

左边是蔗糖，右边两个 \chem{C_6H_{12}O_6} 一个是葡萄糖、一个是果糖——同分异构体，分子式分不出它们，这正是符号的边界。

淀粉的水解（成千上万个葡萄糖单元，用 $n$ 表示）：

\[(\chem{C_6H_{10}O_5})_n + n\chem{H_2O} \rightarrow n\chem{C_6H_{12}O_6}\]

注意淀粉的“单元分子式”比葡萄糖少一个 \chem{H_2O}——每接一个葡萄糖付过一分子水费，账本完全对得上。

而葡萄糖最终当燃料烧掉时——在细胞里叫细胞呼吸（第 52 章起细讲），在火里叫燃烧，净账相同：

\[\chem{C_6H_{12}O_6} + 6\chem{O_2} \rightarrow 6\chem{CO_2} + 6\chem{H_2O}\]

它放能的原因要按第 27 章的规矩说：不是“糖分子里藏着能量”，而是生成 \chem{C=O} 键和 \chem{O-H} 键放出的能量，超过了拆开旧键吸收的能量——断键吸能，成键放能，净账为正。也正因为放能不等于自己会发生（第 28 章），糖在空气里安安静静，需要火或者细胞里的酶机器来启动。

\begin{shiyan}
\textbf{实验一：嚼出来的甜味。}材料：一小勺白米饭（或一小块白面馒头，不加糖）。步骤：放进嘴里，忍住不咽，慢慢嚼满一分钟，留意舌根处的味道变化。观察点：什么时候开始尝到甜？这就是唾液淀粉酶把淀粉水解成麦芽糖。想一想：甜味为什么不是立刻出现的？

\textbf{实验二：碘伏显色。}材料：药店买的碘伏（皮肤消毒剂）、一小勺米饭或一片土豆、一粒方糖、一小片白纸、几滴清水。步骤：把碘伏各滴一滴在米饭、方糖、白纸上，等半分钟看颜色。观察点：米饭和土豆会变成蓝黑色——碘钻进了淀粉螺旋的管道里，和螺旋一起形成一种深色的复合物；方糖几乎不变色；白纸可能微微变蓝，因为很多纸在生产时加了淀粉做“浆料”——这本身就是一个发现。对比三者的颜色深浅，你还能半定量地猜出哪种材料淀粉多。

安全提示：碘伏只许外用，不要入口、不要进眼睛，沾到衣服会染色；滴过碘伏的食物不要再吃。不要加热碘伏，也不要用其他化学品替代它做显色“升级”。
\end{shiyan}

\section{糖在身体里的三份工作}

现在可以清点糖的岗位了。在生命这套化学系统里（第 45 章），糖打着至少三份工。

\textbf{工作一：燃料。}血糖那“一茶匙”葡萄糖是随取随用的零花钱，大脑尤其依赖它。零花钱不够用时，身体打开仓库：肝脏和肌肉里存着几百克糖原。为什么仓库要建成“灌木状”、分支那么密？因为拆糖原的酶只能从每一根侧枝的末端开工，分支越多，末端越多，能同时开工的“拆卸口”就越多——需要快速供能时，几十上百个末端一起放粮。这个设计的数学，章末挑战层再算。至于葡萄糖被拆开后能量怎样一步一步收进“零钱罐”，第 52 章专门讲。

\textbf{工作二：材料。}植物把葡萄糖接成纤维素，做成细胞壁——支撑起从杂草到大树的一切绿色。棉花约九成是纤维素，麻、木材、纸张都是。虾蟹的硬壳、蘑菇的细胞壁用的是几丁质，它的单体是葡萄糖的“带氮改装版”，同样是糖苷键连成的长链。就连你关节里滑溜溜的润滑液、皮肤里填充的凝胶，主要成分也是带负电的长链糖。\textbf{糖是生物界产量最大的建筑材料}：地球上每年新生产的纤维素以千亿吨计，是有机物里的第一大户。

\textbf{工作三：身份标签。}这是最容易被忽略、却最精妙的一份工。每个细胞的表面都挂着许多短糖链——有的接在蛋白质上，有的接在脂质上——像挂在胸前的工牌，供别的细胞、抗体、病毒来“读”。你的血型就是一张糖做的工牌：红细胞表面有一种糖链，末端差一个糖就分出 A、B、O——A 型血在链末端多接一个 N-乙酰半乳糖胺，B 型多接一个半乳糖，O 型什么都不接。一个糖分子的有无，就是免疫系统眼中“自己人”和“入侵者”的分界线：给 A 型的人输进 B 型血，他的抗体会把外来的红细胞当成敌人围攻，这就是输血必须验血型的分子原因。所以严格地说，“血型”其实是“糖型”。细胞靠糖链互相认出对方、精子和卵子靠糖链对上号、流感病毒靠认出细胞表面的特定糖链来找到入侵的门口——这些故事等第 51 章讲细胞膜时再继续。

\section{当糖“出故障”：乳糖、血糖和糖化}

三份工作对应着三类常见的生活问题，正好拿来检验本章的模型。

\textbf{乳糖不耐受。}牛奶里的乳糖是 $\beta$ 接头的双糖，小肠里的乳糖酶负责拆它。几乎所有婴儿都有充足的乳糖酶——否则没法喝奶长大；但断奶之后，很多人的乳糖酶活性会自然下降，这在东亚人群中尤其普遍。没拆开的乳糖一路进到大肠，成了肠道细菌的盛宴：细菌发酵它，产生气体和酸，于是腹胀、肠鸣、拉肚子。注意两点：第一，这是酶多少的差别，不是“对牛奶过敏”（过敏是免疫系统的事，另一套机制）；第二，它有剂量问题——酶少不等于没有，少量多次喝奶、或改喝酸奶（乳酸菌已经帮忙拆掉了不少乳糖），很多人就没事。从进化上看，成年后保留乳糖酶反而是部分人群里的“新变异”——这个故事第 83 章还会回来。

\textbf{血糖。}前面算过，血糖只有一茶匙，高了低了都麻烦。身体靠激素把它维持在一个窄窗口里，其中胰岛素负责“存粮”：让细胞吸收葡萄糖、把多余的接成糖原。当这套调节失灵——或胰岛素分泌不足，或细胞对它的信号反应迟钝——血糖长期偏高，就是糖尿病。这是医学问题，诊断和治疗要交给医生和检验标准；这里只把分子层面的来龙去脉讲清楚。

\textbf{糖化。}还记得水里那不到千分之一的“开链”葡萄糖吗？它链头的醛基很活泼，会跟蛋白质上的氨基自发地粘上去，慢慢发生一连串反应——不用酶，纯化学。你烤面包时表皮变焦糖色、煎肉变褐，就是同一族反应在高温下的快速版本（厨师叫它美拉德反应）。在体温下它慢得多，但从不停工：红细胞里的血红蛋白活多久，就被葡萄糖慢慢“糖化”多久。糖化血红蛋白（HbA1c）占的比例，因此忠实记录了红细胞过去约三个月里泡过的平均血糖浓度——医生用它评估长期血糖控制，比单次扎手指准得多。长期高血糖会加速糖化，生成一类叫“晚期糖基化终末产物”（AGEs）的东西，与血管、眼、肾的慢性损伤有关，这是糖尿病并发症的机制之一。

但这里必须按证据的规矩把话说全：在正常血糖下，糖化是一个缓慢、且身体有清理机制的生理过程，人人都在发生。市面上把“高血糖的病理”偷换成“吃一口糖就毁容”的抗糖产品，证据等级很低。糖不是毒药——它是你的主要燃料；问题出在剂量、代谢状态和个体差异上。这本书不给任何人开饮食处方，但给一把尺子：凡是把“分子机制”直接剪接成“消费恐吓”的说法，都值得多问一句证据在哪里。

\begin{zhengju}
19 世纪末，化学家已经知道葡萄糖是 \chem{C_6H_{12}O_6}、有 6 个碳、能还原银氨溶液。但四个手性碳上羟基的朝向——16 种排布里它到底是哪一种——没人知道。那时连 X 射线晶体学都还没发明（1912 年才有），没有任何仪器能“看见”一个分子的立体结构。

德国化学家费歇尔（Emil Fischer）用的是纯化学逻辑：改造分子，看产物的对称性。他的关键一招是用硝酸把葡萄糖分子的两头（醛基和末端的羟基）都氧化成羧酸。如果某种排布让改造后的分子左右对称，它就失去了旋光性；不对称就还有。再搭配上把糖链加长一节的“升级反应”、把糖链截短一节的“降级反应”，观察每一步产物旋光性的有无和方向——一条条排除，像在 16 个嫌疑人里逐一核对不在场证明。到 1891 年，他把葡萄糖的相对构型定了下来：16 种排布里唯一与所有实验吻合的那一种。1902 年，他拿到诺贝尔化学奖。

但有一件事费歇尔自己承认办不到：他确定的是分子内部各个碳之间的相对关系，至于整套结构的真实版本是他画的那样、还是它的镜像，当时的任何实验都分不出来。他只能“押注”：规定其中一种为标准（这就是 D/L 记号的来源），并清楚地写明这是人为约定。押对的概率是多少？二分之一。

六十年后，1951 年，荷兰晶体学家比吉沃特（Bijvoet）利用 X 射线的“反常散射”效应，第一次测出了分子的绝对构型（他测的是酒石酸盐）。结论：费歇尔当年押对了。整套 D 型糖的立体结构从此有了锚点。

这个故事值得记住三层：第一，看不见的东西可以用排除法推出来，前提是每条证据都能排除掉某些可能；第二，科学家可以“押注”，但必须标明哪里是押注、哪里是证据；第三，押注要等得起——费歇尔的有生之年没等到裁决，这不妨碍他的推理成为教科书。顺带回答前面留的问题：生命为什么几乎只用 D 型糖？一旦最初的酶碰巧按 D 型糖的形状打造，整个代谢网络就被“锁定”了，换手的代价是重建所有酶。可最初那个“碰巧”为什么偏向这一边、还是纯粹随机——老实说，还不知道，这是生命起源研究里仍然开放的谜。
\end{zhengju}

\section{这张糖地图，哪里会失灵}

本章的模型很好用，但有四条边界要交代清楚。

\begin{itemize}[itemsep=2pt]
\item \textbf{六元环不是平的。}哈沃斯投影画成平整的六边形是为了好看。真实的环折成“椅子”形（第 36 章见过环己烷的椅式），羟基有“平伏”和“直立”两种站位，稳定性差别不小——$\beta$ 葡萄糖在水里占多数，正跟它的羟基全占“平伏位”有关。
\item \textbf{占极少数的不等于不重要。}开链葡萄糖不到千分之一，但糖化、变旋、与许多试剂的反应都靠它。环和链在水中不停地互相转化，反应从链那头“抽走”一点，平衡就补上一点——第 31 章的平衡移动原理在这里默默工作。
\item \textbf{本章只讲了葡萄糖家族。}自然界的单糖有成百上千种，糖链可以分支、接硫酸基、接乙酰基，组合出的“糖链语言”复杂程度不亚于蛋白质和核酸。血型只是这门语言的入门词汇。
\item \textbf{血糖是统计意义上的窗口，不是判决书。}单次读数受进食、运动、压力影响；人群里有个体差异。书里的数字用来理解机制，不用来自我诊断——诊断交给重复测量和医学标准。
\end{itemize}

\begin{wuqu}
“糖都是甜的，甜的都叫糖。”——两个方向都错。

先看“糖都是甜的”：淀粉和纤维素是货真价实的糖（多糖），一点也不甜。甜味需要小分子以合适的形状贴进舌头上的甜味受体，几千个葡萄糖连成的大分子根本贴不进去。

再看“甜的都是糖”：阿斯巴甜是两个氨基酸接成的二肽，甜度约是蔗糖的 200 倍，分子骨架跟糖毫无关系；糖精、三氯蔗糖也一样，只是形状碰巧能骗过甜味受体。还有一个苦涩的历史反例：古罗马人用铅器熬葡萄汁，熬出一种甜丝丝的糖浆给酒调味——那甜味来自醋酸铅，俗称“铅糖”，剧毒。甜，但它既不是糖，更不安全。

甜不甜是舌头受体的事，是不是糖是分子结构的事，安不安全是剂量和毒理的事——三件事，三套判据，谁也别代替谁。
\end{wuqu}

\section{回到开场那四样}

现在可以对答案了。

方糖：蔗糖，一个葡萄糖和一个果糖牵手成的双糖。它溶于水靠羟基，加热变褐是糖在高温下分解、重排出的几百种新分子——焦糖香是一场小型化学风暴。

白米饭：淀粉，$\alpha$-1,4 接头盘成的螺旋弹簧。你的淀粉酶认得这个接头，所以嚼着嚼着就甜；碘伏滴上去变蓝黑，是碘钻进了螺旋的管道。

纸巾：纤维素，$\beta$-1,4 接头绷直的长链，链间氢键捆成纤维束。你的身体没有开这把锁的钥匙，所以纸是材料，不是粮食——哪怕它和米饭是同一种积木搭的。

血型：红细胞表面的糖链，末端差一个糖，A、B、O 就此分家。验血型的试纸认的不是“血”，是糖链的末端长什么样。

四样东西，一种积木。差别从来不在砖，在图纸——在接头的朝向、链的形状、以及挂在谁的表面。

\section{糖的下一站}

先补一笔工业和生态的账。纤维素是地球上产量最大的有机物，人类拿它造纸、织布、盖房子；把淀粉和糖发酵成乙醇，又能当燃料——巴西的汽车大量烧甘蔗乙醇。而自然界最大的“纤维素回收队”是真菌和白蚁肚子里的微生物：没有它们拆解 $\beta$ 接头，碳就会在枯枝落叶里积压成山。糖的生产端——植物怎样用阳光把二氧化碳和水接成葡萄糖——要等第 54 章光合作用细讲，那里有一句必须记住的话：氧气来自水，不是来自二氧化碳。

然后埋一个钩子。前面算了：全身血糖只有 5 克，却随时在被几十万亿个细胞抽走当燃料。细胞可不敢“点燃”葡萄糖——燃烧的高温会把自己烧坏。那它怎么把 2870 kJ 的能量安全地榨出来？答案是：把拆糖的过程切成十步，每步只收一小笔，存进一种叫 ATP 的分子零钱。这条流水线叫糖酵解，第 52 章见。而在那之前，下一章先看燃料家族的另一支——脂质：为什么同样一公斤，它存的能量是糖的两倍多。

\begin{tiaozhan}
糖原为什么要长成“灌木”而不是“长绳”？（跳过不影响主线。）算一笔并行账：拆糖原的酶只能从侧枝末端一节一节往里拆。假设一个糖原颗粒存了约 5 万个葡萄糖。如果它是一条不分枝的长绳，全颗粒只有 1 个末端，单位时间最多释放 $v$ 个葡萄糖；如果分支让末端变成 100 个，同一时间的释放速度就是 $100v$。分支的代价是：每岔出一根枝，要多付一分子“水费”（接 $\alpha$-1,6 接头），还要动用另一种专门的酶来“修剪”分支点。肝糖原要在你饿的时候快速平准血糖，所以分支很密；植物的淀粉要存一整个季节，不着急，分支就稀。同一条化学键，搭法不同，一件成了活期存款，一件成了定期。你可以再推一步：分支太密会怎样？末端之间挤到酶伸不进去，反而拆不动——真实糖原的分支密度，恰好落在“末端多”和“不拥挤”的折中点上。
\end{tiaozhan}

\section*{练一练}

\begin{sikaoti}
\item 画粒子图：用六边形表示葡萄糖环，在 1 号碳上画一面“小旗”表示羟基（朝下为 $\alpha$、朝上为 $\beta$）。分别画出两个葡萄糖以 $\alpha$-1,4 和 $\beta$-1,4 相连的样子，标出糖苷键，并用一句话说明为什么你的唾液淀粉酶只拆得动前者。图旁注明：这是人为的表示法。
\item 画物质流图：一粒米饭从入口到入库——淀粉 $\rightarrow$ 麦芽糖 $\rightarrow$ 葡萄糖 $\rightarrow$ 血液 $\rightarrow$ 肝糖原。在每条箭头旁注明：这一步是水解还是缩合？需要谁催化？水分子是进场还是离场？
\item 算能量账：一瓶 500 毫升的可乐含糖约 \ud{53}{g}（按葡萄糖计）。若全部被身体氧化，放能约多少千焦？占一天基础代谢 \ud{7000}{kJ} 的几分之几？（葡萄糖摩尔质量 \ud{180}{g\cdot mol^{-1}}，氧化放能约 \ud{2870}{kJ\cdot mol^{-1}}。）
\item 用“钥匙与接头”的语言解释：为什么乳糖不耐受的人喝牛奶容易腹胀，喝酸奶却往往没事？为什么说这是“正常的遗传差异”而不是一种“病”？从“剂量”角度给出一个可行的喝奶策略。
\item 画信息流图：红细胞表面糖链末端的一个糖（信息）$\rightarrow$ 抗体识别（读取）$\rightarrow$ 凝集反应（后果）。用这幅图解释：为什么 O 型血的人可以给 A 型输血（少量、紧急情况下），反过来却危险？信息差在哪一个分子上？
\item 设计一个对照实验：只用碘伏、清水和厨房材料，检验“某种饼干是否含淀粉”。写出你的实验组、至少两个对照组（一个已知含淀粉，一个已知不含），以及一种可能导致假阳性的干扰（提示：想想纸张实验里的浆料）和排除它的办法。
\end{sikaoti}
